\sequence{Calcul algébrique, équations et inéquations}{Calcul algébrique}{4--20 janvier 2027}
\section{Parenthèses et distributivité}
\begin{methode}[Supprimer des parenthèses]
Après un signe $+$, supprimer les parenthèses conserve les signes des termes. Après un signe $-$, chaque terme à l'intérieur change de signe. Cette règle concerne une somme ou une différence de termes.
\end{methode}
\ex{Supprimer les parenthèses et réduire : $A=x+(2+3x)$ ; $B=x-(2+3x)$ ; $C=-2x+(-4+5x)$ ; $D=-2x-(-4+5x)$.}
\reponse[3.2]{$A=4x+2$ ; $B=-2x-2$ ; $C=3x-4$ ; $D=-7x+4$.}
\begin{cours}[Distributivité]
Pour tous réels $a,b,c,d$ :
\[b(c+d)=bc+bd,\qquad(a+b)(c+d)=ac+ad+bc+bd.\]
\end{cours}
\ex{Développer et réduire :
\begin{align*}
A&=-4(8x+7)+(-8x+5)(5x+9),\\
B&=(6x+9)(-x+7)+(3x+5)(-9x+6),\\
C&=-9(-x-7)-(-x+3)(-8x-8),\\
D&=-5(-6x-7)+(x-7)(5x+8).
\end{align*}}
\reponse[5]{On développe chaque produit, puis on regroupe les termes de même degré :
\begin{align*}
A&=-32x-28-40x^2-72x+25x+45=-40x^2-79x+17,\\
B&=-6x^2+33x+63-27x^2-27x+30=-33x^2+6x+93,\\
C&=9x+63-(8x^2-16x-24)=-8x^2+25x+87,\\
D&=30x+35+5x^2-27x-56=5x^2+3x-21.
\end{align*}}

\fiche{Factoriser avec un facteur commun}
\begin{cours}[Propriété]
$ka+kb=k(a+b)$ et $ka-kb=k(a-b)$. Factoriser transforme une somme en un produit ; développer fait l'opération inverse.
\end{cours}
\ex{Factoriser :
\begin{align*}
A&=15x^2-21x,\\B&=3(2+3x)-(5+2x)(2+3x),\\
C&=(4x-1)(x+6)+(4x-1),\\D&=(1-6x)^2-(1-6x)(2+5x).
\end{align*}}
\reponse[6]{
$A=3x(5x-7)$.\\
$B=(2+3x)[3-(5+2x)]=(2+3x)(-2-2x)=-2(2+3x)(1+x)$.\\
$C=(4x-1)[(x+6)+1]=(4x-1)(x+7)$.\\
$D=(1-6x)[(1-6x)-(2+5x)]=(1-6x)(-1-11x)=(6x-1)(11x+1)$.
Vérifier les facteurs en redéveloppant.}
\ex{Choisir une forme utile de $f(x)=(x-2)(x+3)$ pour calculer $f(2)$, $f(0)$ et chercher ses zéros.}
\reponse[3.2]{La forme factorisée donne immédiatement $f(2)=0$ et les zéros $2,-3$. En $0$, $f(0)=(-2)\times3=-6$ ; la forme développée $x^2+x-6$ donne aussi ce résultat. Le choix dépend de la question.}

\fiche{Identités remarquables : développer}
\begin{cours}[Identités]
\[(a+b)^2=a^2+2ab+b^2,\qquad(a-b)^2=a^2-2ab+b^2,\]
\[(a+b)(a-b)=a^2-b^2.\]
\end{cours}
\preuve[3.2]{Développer $(a+b)(a+b)$ : $a^2+ab+ba+b^2=a^2+2ab+b^2$. Remplacer $b$ par $-b$ donne la deuxième identité. Enfin $(a+b)(a-b)=a^2-ab+ab-b^2=a^2-b^2$.}
\ex{Développer $(x+3)^2$, $(3x-4)^2$, $(x-3)(x+3)$, $(2x+3)(2x-3)$, $(4-3x)(3x+4)$.}
\reponse[3.2]{$x^2+6x+9$ ; $9x^2-24x+16$ ; $x^2-9$ ; $4x^2-9$ ; $16-9x^2$. Dans le dernier produit, réordonner $3x+4$ en $4+3x$.}
\ex{Développer et réduire $A=(2x-3)^2+(x+5)(3-x)$ et $B=2(x+3)+(2x+3)(2x-3)$.}
\reponse[4]{$A=4x^2-12x+9-x^2-2x+15=3x^2-14x+24$. $B=2x+6+4x^2-9=4x^2+2x-3$.}

\fiche{Identités remarquables : factoriser}
\begin{methode}
Repérer les deux carrés ; pour les formes à trois termes, vérifier le double produit $2ab$. Pour une différence de deux carrés, utiliser $a^2-b^2=(a-b)(a+b)$.
\end{methode}
\ex{Factoriser :
\[A=x^2-2x+1,\quad B=25+16x^2-40x,\quad C=4x^2+12x+9,\]
\[D=1-49x^2,\quad E=9x^2-4,\quad F=(2x+3)^2-64.\]}
\reponse[6]{
$A=(x-1)^2$. $B=(4x-5)^2$. $C=(2x+3)^2$.\\
$D=(1-7x)(1+7x)$. $E=(3x-2)(3x+2)$.\\
$F=[(2x+3)-8][(2x+3)+8]=(2x-5)(2x+11)$.
Pour les trois premiers, vérifier le terme intermédiaire ; pour $F$, traiter $2x+3$ comme un seul bloc.}
\ex[Contrôler une identité]{Les égalités $(a+b)^2=a^2+b^2$ et $a^2+b^2=(a-b)(a+b)$ sont-elles vraies pour tous réels ?}
\reponse[3.2]{Non : pour $a=b=1$, la première donnerait $4=2$ et la seconde $2=0$. Le terme manquant dans la première est $2ab$ ; le produit de la seconde vaut $a^2-b^2$.}
\rituel{Calculer $101^2$ mentalement ; donner un directeur de $2x+3y=1$ ; déterminer $[-1;3]\cap[2;5]$.}{$10201$ ; $(-3;2)$ ; $[2;3]$.}

\fiche{Puissances entières relatives}
\begin{cours}[Définition et règles]
Pour $n\geq1$, $a^n$ est le produit de $n$ facteurs égaux à $a$. Pour $a\ne0$, $a^0=1$ et $a^{-n}=1/a^n$. Dans ce cours, $0^0$ n'est pas défini.
Pour $a,b\ne0$ et des entiers relatifs $m,n$ :
\[a^ma^n=a^{m+n},\quad\frac{a^m}{a^n}=a^{m-n},\quad(a^m)^n=a^{mn},\]
\[(ab)^n=a^nb^n,\qquad\left(\frac ab\right)^n=\frac{a^n}{b^n}.\]
Pour des exposants strictement positifs, certaines bases peuvent être nulles si les expressions restent définies.
\end{cours}
\ex{Simplifier $A=2^4\times2^3$, $B=(-5)^{12}/(-5)^3$, $C=(7^3)^6\times7^{-5}$, $D=7^{24}\times3^2\times2^3$. Calculer aussi $(-3)^5$, $5^1$ et $99999^0$.}
\reponse[4]{$A=2^7=128$. $B=(-5)^9=-1953125$. $C=7^{18-5}=7^{13}$. $D=72\times7^{24}$ : on ne peut pas additionner les exposants de bases différentes. Les trois dernières valeurs sont $-243$, $5$, $1$.}
\vigilance{La définition comme produit répété ne vaut pas pour un exposant négatif. $-3^2=-9$, tandis que $(-3)^2=9$.}
\ex{Résoudre mentalement $2^n>100$ pour le plus petit entier naturel $n$. Écrire un programme qui le trouve.}
\reponse[4]{$2^6=64\leq100<128=2^7$, donc $n=7$. Initialiser $n=0$, $p=1$ ; tant que $p\leq100$, remplacer $p$ par $2p$ et $n$ par $n+1$. Le programme est donné en fin de séquence.}

\fiche{Fractions et expressions fractionnaires}
\begin{cours}[Règles et conditions]
Pour les dénominateurs non nuls :
\[\frac ac\pm\frac bc=\frac{a\pm b}{c},\quad\frac ac\pm\frac bd=\frac{ad\pm bc}{cd},\quad\frac ac\times\frac bd=\frac{ab}{cd}.\]
Une fraction $c/d$ possède un inverse seulement si $c\ne0$ et $d\ne0$. Alors
\[\frac ab\div\frac cd=\frac ab\times\frac dc=\frac{ad}{bc}\quad(b,c,d\ne0).\]
\end{cours}
\ex{Calculer et rendre irréductible :
\[A=\left(\frac23+\frac16\right)\times\frac45,\quad B=\frac12+\frac34\times\frac56,\quad C=\frac45\div\frac23\times\frac7{10},\]
\[D=\frac{\left(\frac25+\frac13\right)\times\frac34}{\frac12-\frac16}.\]}
\reponse[5]{$A=\frac56\times\frac45=\frac23$. $B=\frac12+\frac58=\frac98$. $C=\frac45\times\frac32\times\frac7{10}=\frac{21}{25}$. $D=\frac{11}{15}\times\frac34\div\frac13=\frac{11}{20}\times3=\frac{33}{20}$. Dans $C$, multiplication et division sont effectuées de gauche à droite.}
\ex{Simplifier $\dfrac{x^2-1}{x-1}$ en précisant le domaine. Peut-on ensuite prendre $x=1$ ?}
\reponse[3.2]{Pour $x\ne1$, $\frac{(x-1)(x+1)}{x-1}=x+1$. L'expression initiale reste non définie en $1$ ; la simplification ne lève pas la restriction.}

\fiche{Racines carrées et conditions de validité}
\begin{cours}[Définition et propriétés]
Pour $a\geq0$, $\sqrt a$ est l'unique nombre positif ou nul dont le carré vaut $a$. Ainsi $(\sqrt a)^2=a$. Pour tout $x\in\R$, $\sqrt{x^2}=|x|$.
Pour $a,b\geq0$, $\sqrt{ab}=\sqrt a\sqrt b$. Si $b>0$, $\sqrt{a/b}=\sqrt a/\sqrt b$.
\end{cours}
\preuve[3.2]{Le nombre $\sqrt a\sqrt b$ est positif ou nul et son carré vaut $ab$ ; par unicité de la racine carrée positive, il égale $\sqrt{ab}$. Le même argument appliqué au quotient établit la deuxième formule, sous $b>0$. De même $|x|\geq0$ et $|x|^2=x^2$, donc $\sqrt{x^2}=|x|$.}
\ex{Simplifier $\sqrt{18}$ ; $\sqrt2\sqrt8$ ; $\sqrt{7/25}$ ; $\sqrt{27/75}$ ; $\sqrt{(-5)^2}$.}
\reponse[3.2]{$3\sqrt2$ ; $4$ ; $\sqrt7/5$ ; $\sqrt{9/25}=3/5$ ; $5$.}
\vigilance{Il n'existe pas de règle générale $\sqrt{a+b}=\sqrt a+\sqrt b$. Avec $a=4,b=9$, on aurait à tort $\sqrt{13}=5$. Dire que ces expressions sont toujours différentes serait aussi faux, par exemple si $a=0$.}
\ex{Simplifier $2\sqrt{18}-\sqrt{50}+\sqrt8$. Encadrer $\sqrt{20}$ entre deux entiers consécutifs.}
\reponse[3.2]{$6\sqrt2-5\sqrt2+2\sqrt2=3\sqrt2$. Comme $16<20<25$, $4<\sqrt{20}<5$.}

\fiche{Équations du premier degré}
\begin{cours}[Définition]
Résoudre une équation, c'est trouver toutes les valeurs de l'inconnue qui vérifient l'égalité. Deux équations sont équivalentes si elles ont les mêmes solutions. Ajouter une même expression ou multiplier par un même nombre non nul conserve l'équivalence.
\end{cours}
\begin{methode}
Développer si nécessaire, regrouper les termes en $x$, puis isoler $x$. Si son coefficient est nul, examiner l'égalité restante : elle peut être toujours vraie ou impossible. Vérifier dans l'équation de départ.
\end{methode}
\ex{Résoudre :
\begin{align*}
a)&\ 5x-3=2x+9,& b)&\ 3(2x-1)=2(x+4)+7,\\
c)&\ -2(x+3)=4x-(x-1),&d)&\ \frac{x+2}{3}=x-1,\\
e)&\ \frac{-5x+5}{8}=-\frac34,&f)&\ \frac{-3x-1}{9}=-\frac74.
\end{align*}}
\reponse[6]{a) $3x=12$, donc $x=4$. b) $6x-3=2x+15$, donc $x=9/2$. c) $-2x-6=3x+1$, donc $x=-7/5$. d) $x+2=3x-3$, donc $x=5/2$. e) $-5x+5=-6$, donc $x=11/5$. f) $4(-3x-1)=-63$, donc $-12x=-59$ et $x=59/12$.}
\ex{Résoudre $2(x+1)=2x+2$, puis $2(x+1)=2x+3$.}
\reponse[2.4]{La première se réduit à $2=2$ : $S=\R$. La seconde donne $2=3$ : $S=\varnothing$. On ne divise pas par un coefficient nul.}

\fiche{Produit nul et équations avec un quotient}
\begin{cours}[Propriétés]
$AB=0\iff A=0$ ou $B=0$. Pour $B\ne0$, $A/B=0\iff A=0$. Pour résoudre $A(x)/B(x)=k$, imposer d'abord $B(x)\ne0$, puis résoudre $A(x)=kB(x)$ et écarter les valeurs interdites.
\end{cours}
\ex{Résoudre $(4x+6)(3-7x)=0$ ; $3x(-7x+5)=0$ ; $\dfrac{2x+1}{3x-5}=0$.}
\reponse[3.2]{Solutions : $\{-3/2;3/7\}$ ; $\{0;5/7\}$ ; $\{-1/2\}$, avec valeur interdite $5/3$ dans le dernier cas.}
\ex{Factoriser puis résoudre $(1-6x)^2-(1-6x)(2+5x)=0$ et $(4x-1)(x+6)+(4x-1)=0$.}
\reponse[3.2]{Première : $(6x-1)(11x+1)=0$, donc $S=\{1/6;-1/11\}$. Deuxième : $(4x-1)(x+7)=0$, donc $S=\{1/4;-7\}$.}
\ex{Résoudre $\dfrac{x+1}{x-2}=3$, puis $\dfrac{x^2-4}{x-2}=4$.}
\reponse[4]{Dans les deux cas, $x\ne2$. Première : $x+1=3x-6$, donc $x=7/2$, autorisé. Deuxième : pour $x\ne2$, le quotient vaut $x+2$ ; l'équation impose $x=2$, interdit. Donc aucune solution.}

\fiche{Équations avec un carré ; isoler une variable}
\begin{cours}[Rappel]
Dans $\R$, $x^2=a$ a deux solutions $\pm\sqrt a$ si $a>0$, une solution $0$ si $a=0$, et aucune si $a<0$.
\end{cours}
\ex{Résoudre $x^2=0$ ; $x^2=-8$ ; $x^2=7$ ; $25x^2=81$ ; $(x+2)^2=9$.}
\reponse[4]{Dans l'ordre : $\{0\}$ ; $\varnothing$ ; $\{-\sqrt7;\sqrt7\}$ ; $\{-9/5;9/5\}$ ; $\{-5;1\}$. Pour la dernière, $x+2=-3$ ou $x+2=3$.}
\begin{methode}[Isoler une variable dans une formule]
Décrire les opérations effectuées sur les deux membres ; préciser les conditions qui permettent de diviser ou de prendre une racine carrée.
\end{methode}
\ex{Isoler $h$ dans $V=\pi r^2h$ avec $r>0$ ; isoler $v$ dans $E=\frac12mv^2$ avec $m>0$, $E\geq0$ et $v\geq0$ ; isoler $x$ dans $y=ax+b$ avec $a\ne0$.}
\reponse[4]{$h=V/(\pi r^2)$. Puis $v^2=2E/m$ ; comme $v\geq0$, $v=\sqrt{2E/m}$. Enfin $x=(y-b)/a$. Sans la condition $v\geq0$, deux signes seraient possibles pour $E>0$.}
\rituel{Donner le domaine de $1/(x+3)$ ; calculer $\sqrt{49}$ ; la pente entre $(0;2)$ et $(3;8)$.}{$\R\setminus\{-3\}$ ; $7$ ; $2$.}

\fiche{Inégalités et comparaison}
\begin{cours}[Règles]
On peut additionner des inégalités de même sens : $a\leq b$, $c\leq d\Rightarrow a+c\leq b+d$. Ajouter un même réel conserve le sens. Multiplier ou diviser par un réel strictement positif conserve le sens ; par un réel strictement négatif, le renverse.
Comparer $A$ et $B$ revient à étudier le signe de $A-B$. Si $A,B>0$, on peut aussi comparer $A/B$ à $1$.
\end{cours}
\ex{1. Si $x\leq-4$, que dire de $3x-4$ ? 2. Si $x>-6$, que dire de $2x-3$ ? 3. Si $7\leq x\leq16$ et $3\leq y\leq12$, encadrer $8x+7y$. 4. Si $-10\leq x\leq-6$, encadrer $2x-6$.}
\reponse[4]{1. $3x-4\leq-16$. 2. $2x-3>-15$. 3. $56\leq8x\leq128$ et $21\leq7y\leq84$, donc $77\leq8x+7y\leq212$. Les bornes sont atteignables. 4. $-26\leq2x-6\leq-18$.}
\ex{Résoudre $3x+1\geq x-7$ ; $2x-11\leq4x-5$ ; $6x+19\leq2x-3$.}
\reponse[4]{1. $2x\geq-8$, soit $[-4;+\infty[$. 2. $-6\leq2x$, soit $[-3;+\infty[$. 3. $4x\leq-22$, soit $]-\infty;-11/2]$.}

\fiche{Signe d'une expression affine et d'un produit}
\begin{cours}[Propriété]
Pour $a\ne0$, $ax+b$ s'annule en $x_0=-b/a$ ; son signe à droite est celui de $a$, à gauche l'opposé. Si $a=0$, le signe est celui de la constante $b$. Pour un produit, multiplier les signes des facteurs.
\end{cours}
\ex{Dresser les tableaux de signes de $3-x$, $2x-6$, puis $P(x)=(2x-6)(x+1)$.}
\reponse[5]{Les deux expressions affines s'annulent en $3$ : $3-x$ est positif à gauche, négatif à droite ; $2x-6$ a les signes inverses.
\begin{center}\renewcommand{\arraystretch}{1.5}\begin{tabular}{c|ccccccc}
$x$&$-\infty$&&$-1$&&$3$&&$+\infty$\\\hline
$2x-6$&&$-$&$-$&$-$&$0$&$+$&\\
$x+1$&&$-$&$0$&$+$&$+$&$+$&\\\hline
$P(x)$&&$+$&$0$&$-$&$0$&$+$&
\end{tabular}\end{center}}
\ex{Résoudre $(7+3x)(2x+6)>0$.}
\reponse[5]{Les zéros sont $-7/3$ et $-3$, à ranger $-3<-7/3$. Le produit est positif à l'extérieur de ces racines et négatif entre elles. L'inégalité est stricte : $S=]-\infty;-3[\cup]-7/3;+\infty[$.}

\fiche{Signe d'un quotient et synthèse algorithmique}
\begin{methode}
Chercher les valeurs interdites, puis les zéros du numérateur. Ranger toutes les valeurs critiques et appliquer la règle des signes. Sur la ligne du quotient, une double barre indique une valeur interdite.
\end{methode}
\ex{Dresser le tableau de signes de $Q(x)=\dfrac{-4x+6}{x-3}$, puis résoudre $\dfrac{2x-1}{-4x+3}\leq0$.}
\reponse[5]{Pour $Q$, zéro $3/2$, interdit $3$. Signes : négatif sur $]-\infty;3/2[$, nul en $3/2$, positif sur $]3/2;3[$, non défini en $3$, négatif sur $]3;+\infty[$.
Pour la seconde expression, zéro $1/2$, interdit $3/4$. Les signes sont $-,+,-$ sur les trois intervalles. Ainsi $S=]-\infty;1/2]\cup]3/4;+\infty[$.}
\begin{lstlisting}
def premiere_puissance(seuil):
    # seuil reel >= 1
    n, p = 0, 1
    while p <= seuil:
        n = n+1
        p = 2*p
    return n, p
\end{lstlisting}
\ex{Que renvoie \texttt{premiere\_puissance(100)} ? Quel lien reste vrai entre \texttt{p} et \texttt{n} à chaque étape ? Comment tester ensuite le signe d'un nombre en Python ?}
\reponse[3.2]{Le résultat est $(7,128)$. L'invariant est $p=2^n$, vrai à l'initialisation puis conservé. Pour le signe, tester d'abord \texttt{x > 0}, puis \texttt{x < 0}, sinon $x=0$. Ne pas remplacer l'inégalité large de la boucle par une stricte sans modifier le problème posé.}
